% ZOMBI2 manual — LaTeX preamble, injected via pandoc --include-in-header.
% Kept minimal and re-load-safe (pandoc's default template already loads the core packages).

\usepackage{graphicx}
\setkeys{Gin}{width=0.6\linewidth,keepaspectratio}   % figures default to 60% text width

\usepackage{booktabs}
\usepackage{microtype}

% --- Glyph coverage ---------------------------------------------------------
% XeLaTeX drops a character it has no glyph for *silently* — it simply vanishes
% from the PDF, leaving only a "Missing character" line in the build log. Latin
% Modern lacks several characters the chapters use, so both places they appear
% need covering, and they need covering differently.
%
% In verbatim (code blocks, the ASCII diagrams in Ch5) catcodes are frozen and no
% macro can intervene, so the only fix is a font that genuinely has the glyphs.
% DejaVu Sans Mono covers box drawing, Greek and the maths operators. Loaded by
% filename so kpathsea finds it inside TeX Live, with no dependency on the
% system font database (which on a CI runner has nothing in it).
\setmonofont{DejaVuSansMono.ttf}[
  BoldFont       = DejaVuSansMono-Bold.ttf,
  ItalicFont     = DejaVuSansMono-Oblique.ttf,
  BoldItalicFont = DejaVuSansMono-BoldOblique.ttf,
  Scale          = MatchLowercase]

% In running text the fix is a mapping to the LaTeX equivalent, which keeps the
% body font as it is. Only characters the book actually uses are listed; to check
% for new gaps, grep the build log for "Missing character".
\usepackage{newunicodechar}
\newunicodechar{σ}{\ensuremath{\sigma}}
\newunicodechar{⁻}{\textsuperscript{\textminus}}

% --- Long file names in tables ----------------------------------------------
% Appendix B's tables carry names like `sequences_species_phylogram_complete.nwk`,
% which LaTeX treats as one unbreakable box and pushes out past the right margin.
% Let a line break after an underscore instead. \allowbreak, not hyphenation:
% a hyphen inserted into a file name would be read as part of the name, which in
% a reference table of file names is simply wrong. Pandoc writes an underscore in
% a code span as \_, so redefining that covers every one of them.
\let\zombitextunderscore\_
\renewcommand{\_}{\zombitextunderscore\allowbreak}

% Reference tables are dense and set narrow; a size down keeps them inside the
% text block and is the usual typography for a table of this kind.
\usepackage{etoolbox}
\AtBeginEnvironment{longtable}{\small}

% Compact chapter headings — drop the ~1in of blank space the book class puts above them.
\usepackage{titlesec}
\titleformat{\chapter}[display]{\normalfont\huge\bfseries}%
  {\chaptertitlename\ \thechapter}{6pt}{\Huge}
\titlespacing*{\chapter}{0pt}{0pt}{18pt}

% Pin figures in place: they render exactly where they appear in the source, instead of
% floating forward into a later section (the float package's [H] placement, set as the
% default for every figure). \intextsep still controls the space around them.
\usepackage{float}
\floatplacement{figure}{H}
\setlength{\intextsep}{10pt plus 2pt minus 2pt}

\usepackage{xcolor}
\definecolor{zteal}{HTML}{0B6B63}
\definecolor{zwarn}{HTML}{B03A2E}
\definecolor{ztip}{HTML}{1E8449}

% Callout boxes (see callouts.lua: ::: note / ::: warning / ::: tip)
\usepackage[most]{tcolorbox}
\newtcolorbox{notebox}{breakable,enhanced,colback=zteal!4,colframe=zteal,
  coltitle=white,fonttitle=\bfseries,title=Note,boxrule=0.6pt,
  left=2.5mm,right=2.5mm,top=1.2mm,bottom=1.2mm}
\newtcolorbox{warnbox}{breakable,enhanced,colback=zwarn!4,colframe=zwarn,
  coltitle=white,fonttitle=\bfseries,title=Warning,boxrule=0.6pt,
  left=2.5mm,right=2.5mm,top=1.2mm,bottom=1.2mm}
\newtcolorbox{tipbox}{breakable,enhanced,colback=ztip!4,colframe=ztip,
  coltitle=white,fonttitle=\bfseries,title=Tip,boxrule=0.6pt,
  left=2.5mm,right=2.5mm,top=1.2mm,bottom=1.2mm}

\usepackage{caption}
\captionsetup{font=small,labelfont=bf,labelsep=period}
